% TOP_PROBLEM: {"id": 3064, "title": "Bases of many weights", "category_id": 2, "category": {"id": 2, "name": "combinatorics", "display_name": "Combinatorics", "description": "Counting problems, graph theory, discrete structures.", "slug": "combinatorics", "order_index": 2, "created_at": "2026-07-31T15:26:25.671Z"}, "status": "open", "problem_number": "OPG-369", "set_id": 12}
% FIRST_POSTED: 2026-10-01
% ENTRY_KIND: statement_only
% UnsolvedMath ID 3064; source catalogue code OPG-369
% !TeX program = lualatex
% Definition and sources only; this file is not an LLM solution attempt.
\documentclass[11pt,a4paper]{article}
\usepackage[margin=25mm]{geometry}
\usepackage{fontspec}
\setmainfont{lmroman10-regular.otf}[BoldFont=lmroman10-bold.otf,
 ItalicFont=lmroman10-italic.otf,BoldItalicFont=lmroman10-bolditalic.otf]
\usepackage{amsmath,amssymb,mathtools}
\usepackage{unicode-math}
\setmathfont{latinmodern-math.otf}
\AtBeginDocument{\let\setminus\smallsetminus}
\usepackage{xurl,hyperref}
\hypersetup{colorlinks=true,urlcolor=blue}
\setcounter{secnumdepth}{0}
\setlength{\parindent}{0pt}
\setlength{\parskip}{5pt}
\setlength{\emergencystretch}{3em}
\begin{document}
\section*{Bases of many weights}\label{3064}
{\small UnsolvedMath ID 3064 · OPG-369 · Combinatorics · OpenGarden}
\subsection{Definitions and mathematical statement}
Let $M=(E,\mathcal I)$ be a matroid on a finite set $E$: $\mathcal I$ is
a nonempty hereditary family of subsets, and if $I,J\in\mathcal I$ with
$|I|<|J|$, some $e\in J\setminus I$ has $I\cup\{e\}\in\mathcal I$.
Write $r(X)=\max\{|I|:I\subseteq X,\ I\in\mathcal I\}$.
A basis is a maximal independent set, necessarily of size $r(E)$.

Let $G$ be an additive abelian group and assign a weight $w(e)\in G$
to each element. Define the set of distinct basis weights and its stabilizer by
\[
 S=\left\{\sum_{e\in B}w(e):B\text{ a basis of }M\right\},\qquad
 H=\{g\in G:g+S=S\}.
\]
The set $S$ is finite and nonempty, and $H$ is a finite subgroup, even
if $G$ is infinite. Let $G/H$ denote the set of its additive cosets.
For a coset $Q$, write $w^{-1}(Q)=\{e\in E:w(e)\in Q\}$.

The conjecture asserts
\[
 |S|\ge |H|\left(1-r(E)+
                 \sum_{Q\in G/H}r(w^{-1}(Q))\right).
\]
Only finitely many summands are nonzero, so this is an ordinary finite
integer bound. The left side counts different sums, not the number of
bases that attain them. Repeated element weights are allowed, and the
rank of each weight class retains the dependence structure of the
matroid. In rank zero the empty basis has weight zero, giving the same
well-defined conventions.
\subsection{Short English statement}
Does the number of distinct weighted matroid-basis sums satisfy the stated lower bound after accounting for its translation stabilizer?
\subsection{Sources}
\begin{itemize}
\item[S1] ULAM.AI and the UnsolvedMath Contributors, supplied problems.json, record 3064 (OPG-369), OpenGarden collection. Catalogue identity and source transcription; CC BY 4.0.
\url{https://huggingface.co/datasets/ulamai/UnsolvedMath}.
\item[S2] Open Problem Garden, Bases of many weights. Original problem and discussion; the mathematical formulation below is expanded from the supplied source record.
\url{http://www.openproblemgarden.org/op/bases_of_many_weights}.
\end{itemize}
\end{document}
